% Values: PRECISE, FAITHFUL, ILLUMINATING, ECONOMICAL, TRACTABLE
% Evidence planned:
% - PRECISE: every effect value (task, command, subscription, process
%   identifier, program value, JSON value, boundary value) defined by set
%   constructions;
%   the result of APPLY on every basic value of Task, Process, Cmd, Sub and
%   Platform that is not written in Elm, in one figure (fig:effects-apply);
%   task reduction by rules; Enc and Dec by one rule per admissible type
%   form; the message of every kind of leaf by rules.
% - FAITHFUL: D10 (tasks, commands, subscriptions as values of the
%   language), D11 (what crosses a port, both ways, with a table), D2 and
%   D3 (no exceptions, no run-time errors: a bad incoming value is rejected
%   at the boundary, A71; a failing process ends alone, A79; kill of an
%   ended process does nothing, A80); A60 (Elm-written basis functions
%   defined by their text); A66, A67 (one order convention); A73 to A77,
%   A81, A83, A84, A85 cited where they shape the text; A92 (Bytes at the
%   boundary, as byte sequences that are not JSON values).
% - ILLUMINATING: comments on the silent and labelled steps, the suspension
%   points of a process (A78), the unconstrained process identifier of
%   spawn, the Maybe (Maybe t) and Maybe () collapse, the extra tuple
%   elements and record fields that decoding ignores, unit at ports versus
%   flags, and the outgoing port leaf that gives no message.
% - TRACTABLE: stated properties: the form of the steps of a task; Enc and
%   Dec are partial functions, Enc is defined on every payload; decoding
%   an encoding gives the value back, with the exceptions named.
% - ECONOMICAL: one rule for all primitives through Out; options merge the
%   silent and labelled inner steps and the two tuple arities; none is
%   batch with no components; perform, attempt, map, sequence, Cmd.none
%   and Sub.none get no definition beyond their Elm text.
\chapter{Tasks, Commands, Subscriptions and Ports}\label{ch:effects}

This chapter defines the effect values of Standard Elm: the tasks, which
describe work that may wait, fail and start other processes; the commands and
subscriptions, which a program hands to the platform to have tasks performed
and values exchanged with its host (the platform and the host are as in
Chapter~\ref{ch:introduction}); and the program values. % D10
It defines the boundary values that cross a port, which are JSON values in
which byte sequences may stand, and the encoding and decoding of the values of
every type that a port admits, which settles exactly what can and cannot be
sent over a port. % D11, A92
Its sentences infer that a task reduces to another task, silently or by
performing a primitive action; that a value encodes as a boundary value, or a
boundary value decodes as a value, at a type; and that a leaf of a command or
subscription, receiving an input, gives a message. These objects are
defined here; the running of a program, which creates and schedules processes
and handles their messages one at a time, is defined by
Chapter~\ref{ch:programs} in terms of them.
Every phrase whose value is used here is assumed to have elaborated
(Chapter~\ref{ch:static-modules}). % S4

\section{Simple Objects}\label{sec:effects-simple}

The simple objects of this chapter are shown in
Figure~\ref{fig:effects-simple}; the classes of Chapter~\ref{ch:dynamic-core}
are used as there.

\begin{figure}[htbp]
\centering
\begin{objtable}
\objrow{\pid}{\ProcId}{process identifiers}
\objrow{\pi}{\effPrim}{primitives}
\objrow{\effp\ \text{or}\ \var}{\cls{Var}}{port names}
\objrow{}{\effMgrPlat}{managers of the platform}
\objrow{}{\effJNum}{JSON numbers}
\objrow{}{\effByt}{byte sequences}
\end{objtable}
\caption{Simple Semantic Objects of the Effects}
\label{fig:effects-simple}
\end{figure}

A \emph{process identifier} names a process (Section~\ref{sec:effects-tasks});
ProcId is an infinite set. The class Prim of \emph{primitives} is
\[
  \effPrim=\{\effsleep,\effspawn,\effkill\}\cup\effPrimPlat ,
\]
where $\effsleep$, $\effspawn$ and $\effkill$ are the primitives of the module
\kw{Process}, and $\effPrimPlat$ is a set of further primitives (reading a
clock, an HTTP request, and so on) which is a parameter of the platform on
which a program runs. % A83
Each primitive $\pi$ has an \emph{arity} $\arity(\pi)\geq 1$; the arity of
each of $\effsleep$, $\effspawn$ and $\effkill$ is 1.
A port name $\effp$ is the variable of a \kw{port} declaration
(Chapter~\ref{ch:syntax-modules}). Port names are global to a program and
distinct from one another (Section~\ref{sec:syntax-modules-restrictions}). % A84
$\effMgrPlat$, a parameter of the platform, is a set of names of
\emph{managers}, other than the managers of tasks and ports defined in
Section~\ref{sec:effects-commands}, and disjoint from the port names. % A83, A84

$\effJNum$ is the set of \emph{JSON numbers}: the rational numbers
$m\times 10^{e}$ with $m$ and $e$ integers, that is, the numbers a JSON text
can spell. Every integer of $\IntRange$ and every finite floating-point number
is a member of $\effJNum$; we omit the injections from these numbers into
$\effJNum$. The injection takes both floating-point zeros, $0$ and $-0$, to the JSON
number $0$: a JSON number has no sign apart from its value. % A88
We write $\effStr$ for the set of strings of SVal, $\effFlt$ for the set of
floating-point numbers of SVal, and $\effFltFin\subseteq\effFlt$ for its finite
members, so that $\effFlt\setminus\effFltFin$ holds the two infinities and
NaN. For $j\in\effJNum$, $\efffloat(j)$ is the floating-point number nearest
$j$, rounded as the value of a floating-point constant is
(Section~\ref{sec:dynamic-core-simple}); $\efffloat(0)$ is the float $0$, not
$-0$. % A88

$\effByt$ is the set of \emph{byte sequences}: the finite sequences of
integers from $0$ to $255$. We write $\effbytes(o_1\cdots o_k)$, with $k\geq 0$
and $0\leq o_i\leq 255$ for each $i$, for the byte sequence of
$o_1,\ldots,o_k$ in order. A byte sequence is not a JSON number, a string or an
array: $\effByt$ is disjoint from $\effJNum$, from $\effStr$ and from every
other class of this chapter. % A92

\section{Compound Objects}\label{sec:effects-compound}

The compound objects of this chapter are shown in
Figure~\ref{fig:effects-compound}. The classes EffVal and JVal, which
Chapter~\ref{ch:dynamic-core} includes in Val
(Figure~\ref{fig:dynamic-core-compound}), are defined here. TaskVal, CmdVal,
SubVal and JVal are defined together with the classes of
Figure~\ref{fig:dynamic-core-compound}; all are the smallest classes closed
under the constructions of both figures. The class BVal of boundary values,
which includes JVal, is defined here too; its other members are not in
Val. % A92

\begin{figure}[htbp]
\centering
\begin{objtable}
\objdef{\ev}{\EffVal}{\TaskVal\cup\CmdVal\cup\SubVal\cup\ProcId\cup\ProgVal}
\objdef{\tk}{\TaskVal}{(\{\effsucceed,\efffail\}\times\Val)}
& & & \phantom{{}={}}\cup(\{\effandThen,\effonError\}\times\Val\times\TaskVal) & \\
& & & \phantom{{}={}}\cup\textstyle\bigcup_{k\geq 1}(\{\effprim\}\times\effPrim\times\Val^k) & \\
\objdef{\cm}{\CmdVal}{\textstyle\bigcup_{n\geq 0}(\{\effbatch\}\times\CmdVal^n)}
& & & \phantom{{}={}}\cup(\{\effmap\}\times\Val\times\CmdVal) & \\
& & & \phantom{{}={}}\cup(\{\effleaf\}\times\effMgr\times(\Val\cup\BVal)) & \\
\objdef{\sbv}{\SubVal}{\textstyle\bigcup_{n\geq 0}(\{\effbatch\}\times\SubVal^n)}
& & & \phantom{{}={}}\cup(\{\effmap\}\times\Val\times\SubVal) & \\
& & & \phantom{{}={}}\cup(\{\effleaf\}\times\effMgr\times\Val) & \\
\objdef{\effm}{\effMgr}{\{\effmtask\}\cup(\{\effmout\}\times\cls{Var})}
& & & \phantom{{}={}}\cup(\{\effmin\}\times\cls{Var}\times\Type)\cup\effMgrPlat & \\
\objdef{\effa}{\effAction}{\textstyle\bigcup_{k\geq 1}\effPrim\times\Val^k\times\TaskVal}
\objdef{\pv}{\ProgVal}{\{\effprog\}\times\Record}
\objdef{j}{\BVal}{\{\effnull,\efftrue,\efffalse\}\cup\effJNum\cup\effStr\cup\effByt}
& & & \phantom{{}={}}\cup\textstyle\bigcup_{n\geq 0}\BVal^n\cup(\effStr\finmap\BVal) & \\
\objdef{}{\JVal}{\{\effnull,\efftrue,\efffalse\}\cup\effJNum\cup\effStr}
& & & \phantom{{}={}}\cup\textstyle\bigcup_{n\geq 0}\JVal^n\cup(\effStr\finmap\JVal) & \\
\end{objtable}
\caption{Compound Semantic Objects of the Effects}
\label{fig:effects-compound}
\end{figure}

We write the members of these classes as follows:
\[
  \begin{array}{ll}
  \text{tasks:} & \effsucceed\,v,\quad \efffail\,v,\quad \effandThen(v,\tk),\\
    & \effonError(v,\tk),\quad \effprim(\pi,v_1\cdots v_k);\\
  \text{commands:} & \effbatch(\cm_1\cdots\cm_n),\quad \effmap(v,\cm),\quad
    \effleaf(\effm,v);\\
  \text{managers:} & \effmtask,\quad \effmout\,\effp,\quad \effmin(\effp,\tau);\\
  \text{actions:} & \effdone{\pi(v_1\cdots v_k)}{\tk};\\
  \text{program values:} & \effprog\,r.
  \end{array}
\]
Subscriptions are written as commands are, and an action
$\effdone{\pi(v_1\cdots v_k)}{\tk}$ is the triple $(\pi,v_1\cdots v_k,\tk)$. We write $\effnone$ for
$\effbatch(\,)$, the batch with no components, in CmdVal and in SubVal. In a
primitive task $\effprim(\pi,v_1\cdots v_k)$, always $k=\arity(\pi)$.

A \emph{boundary value} is $\effnull$, $\efftrue$, $\efffalse$, a number, a
string, a byte sequence, an \emph{array}, which is a member
$\effJarr{j_1,\ldots,j_n}$ of $\BVal^n$, or an \emph{object}, which is a
finite map from strings to boundary values. An object has no order on its
keys. % A76
A boundary value is thus a JSON value in which a byte sequence may stand
wherever a value may. % A92
A \emph{JSON value} is a boundary value in which no byte sequence occurs; JVal
is the class of JSON values, and $\JVal\subseteq\BVal$. A byte sequence is
not a JSON value. The metavariable $j$ ranges over BVal; where only a JSON value
is meant we write $j\in\JVal$.
The values of the type \kw{Json.Encode.Value} (also called
\kw{Json.Decode.Value}) are exactly the JSON values, the members of JVal, so
no byte sequence occurs in one. % A74, A92
As with the other classes, JVal is disjoint from the rest of Val, and we omit
the injection: a JSON string and the String value with the same characters are
different values, and $\effnull$ is not the unit value $()$. A boundary value
that is not a JSON value is not a value of the language; such a boundary value
occurs only at a port, and in the leaves of outgoing ports, which CmdVal
admits.

For values $v_1,\ldots,v_n$ we write $[v_1,\ldots,v_n]$ for the list value
$v_1\mathbin{\kw{::}}\cdots\mathbin{\kw{::}}v_n\mathbin{\kw{::}}\kw{[]}$, and
$[\,]$ when $n=0$. We write the value to which the initial dynamic basis
(Appendix~\ref{ch:appendix-dynamic-basis}) binds a long identifier of the
basis by that identifier in typewriter, as in \kw{Task.succeed}; the basic
values \kw{Array.fromList} and \kw{Array.toList} of
Appendix~\ref{ch:appendix-dynamic-basis}
(Section~\ref{sec:appendix-dynamic-basis-arrays}), which no identifier of the
basis spells, are written likewise. Finite multisets are written $\effmset{\effc_1,\ldots,\effc_n}$,
and $\uplus$ is the union of multisets, which adds multiplicities.

\section{Inference Rules}\label{sec:effects-rules}

The semantic rules of this chapter allow sentences of the forms
\[
  \tk\effred{\opt{\effa}}\tk',\qquad
  \Enc{\tau}(v)=j,\qquad
  \Dec{\tau}(j)=v,\qquad
  \effrecv{\effc}{\effx}{v}
\]
to be inferred. The first is pronounced ``$\tk$ reduces to $\tk'$'', and,
when the label $\effa$ is present, ``performing $\effa$''; a reduction without
a label is \emph{silent}. The second is pronounced ``$v$ encodes at $\tau$ as
$j$'', and the third ``$j$ decodes at $\tau$ as $v$''; there $\tau$ is a type
with $\Port(\tau)$ (Section~\ref{sec:static-modules-ports}). In the fourth,
$\effc$ is a command or subscription value with no $\effbatch$ in it (a
\emph{leaf command}, Section~\ref{sec:effects-commands}), the input $\effx$ is a
value or a pair $(\effp,j)$ of a port name and a boundary value, and the sentence
is pronounced ``$\effc$, receiving $\effx$, gives the message $v$''. Premises also use
the sentence $\japp{v_1}{v_2}{v}$ of the application of values
(Section~\ref{sec:dynamic-core-apply}). Some hypotheses in rules are not of
these forms; they are side conditions.

The convention for options is the same as for the Core static semantics. In
particular, in a task reduction whose label is written $\opt{\effa}$, the
label is present on every arrow of the rule or absent from every one. As in
the Core (Section~\ref{sec:dynamic-core-rules}), sentences have no states and
there are no exception packets, so every rule is given in full. % D2
We allow the compound variable $\cm/\sbv$ to range over
$\CmdVal\cup\SubVal$.

$\Enc{\tau}$ and $\Dec{\tau}$ are defined by the rules alone: $\Enc{\tau}(v)$
is the boundary value $j$ for which $\Enc{\tau}(v)=j$ can be inferred, and is
undefined when there is none, and likewise for $\Dec{\tau}(j)$
(Section~\ref{sec:effects-ports}).

\section{Tasks}\label{sec:effects-tasks}

A task value describes work; it does nothing until a process performs it, and
a task value performed twice does its work twice. % D10
The task $\effsucceed\,v$ has \emph{succeeded} with $v$, and $\efffail\,v$ has
\emph{failed} with $v$; these two are the \emph{outcomes}. The task
$\effandThen(v,\tk)$ performs $\tk$ and, if it succeeds with $v'$, goes on
with the task to which $v$ applied to $v'$ evaluates; $\effonError(v,\tk)$ is
the same for failure. A primitive task $\effprim(\pi,v_1\cdots v_k)$ performs
the primitive $\pi$ on $v_1\cdots v_k$.

A \emph{process} is a process identifier together with the task that it is
reducing. A process reduces its task by the rules below, one step at a time,
until the task is an outcome; then the process has \emph{ended}. A silent step
concerns that process alone. A labelled step performs a primitive action
$\effdone{\pi(v_1\cdots v_k)}{\tk}$, whose outcome $\tk$ is a member of the
set
\[
  \effOut(\pi,v_1\cdots v_k)
\]
of possible outcomes, and which may act on the world outside the process. For
the primitives of \kw{Process},
\[
  \begin{array}{rcl}
  \effOut(\effsleep,v) &=& \{\effsucceed\,()\},\\
  \effOut(\effspawn,\tk) &=& \{\effsucceed\,\pid \mid \pid\in\ProcId\},\\
  \effOut(\effkill,\pid) &=& \{\effsucceed\,()\},
  \end{array}
\]
and their actions mean the following.
\begin{itemize}
\item The action $\effdone{\effsleep(v)}{\effsucceed\,()}$, where
  $v\in\effFlt$, takes place no sooner than $\max(0,v)$ milliseconds after the process reached
  $\effprim(\effsleep,v)$, when $v\in\effFltFin$; when $v$ is an infinity or
  NaN, when it takes place is not defined. % A81
\item $\effdone{\effspawn(\tk)}{\effsucceed\,\pid}$ starts a new process,
  named by $\pid$, reducing $\tk$. The identifier $\pid$ names no process
  that has been started before. The new process is independent of the one
  that started it: its outcome is delivered to no other process, and when it
  fails it ends, and nothing else happens. % A79
\item $\effdone{\effkill(\pid)}{\effsucceed\,()}$ ends the process named by
  $\pid$ at once, so that it takes no further step, when that process has not
  ended and is not the process performing the action. Otherwise the action has
  no effect. % A80
  % A90: a process that kills itself is not ended, and continues with
  % succeed (); both implementations behave so (the kill binding resumes its
  % own caller).
\end{itemize}
For a primitive $\pi\in\effPrimPlat$, the set $\effOut(\pi,v_1\cdots v_k)$,
which may hold failures, and the meaning of its actions are given by the
platform. % A83
Chapter~\ref{ch:programs} defines the processes of a running program, and
it is there that the conditions above on when an action may take place, and
on what it does to other processes, are enforced.

\rulegroup{Task Reduction}{\tk\effred{\opt{\effa}}\tk'}

\defnrule[r:e-task-andthen-succeed]
  {\japp{v}{v'}{\tk}}
  {\effandThen(v,\effsucceed\,v')\effred{}\tk}

\defnrule[r:e-task-andthen-fail]
  {}
  {\effandThen(v,\efffail\,v')\effred{}\efffail\,v'}

\defnrule[r:e-task-onerror-fail]
  {\japp{v}{v'}{\tk}}
  {\effonError(v,\efffail\,v')\effred{}\tk}

\defnrule[r:e-task-onerror-succeed]
  {}
  {\effonError(v,\effsucceed\,v')\effred{}\effsucceed\,v'}

\defnrule[r:e-task-andthen-inner]
  {\tk\effred{\opt{\effa}}\tk'}
  {\effandThen(v,\tk)\effred{\opt{\effa}}\effandThen(v,\tk')}

\defnrule[r:e-task-onerror-inner]
  {\tk\effred{\opt{\effa}}\tk'}
  {\effonError(v,\tk)\effred{\opt{\effa}}\effonError(v,\tk')}

\defnrule[r:e-task-prim]
  {\tk\in\effOut(\pi,v_1\cdots v_k)}
  {\effprim(\pi,v_1\cdots v_k)\effred{\effdone{\pi(v_1\cdots v_k)}{\tk}}\tk}

\begin{rulecomments}
\rcomment{r:e-task-andthen-succeed} The task $\tk$ to which $v$ applied to
$v'$ evaluates is not reduced by this step; it is the task the process
performs next. If the application has no result (for instance, it reaches
\kw{Debug.todo}), the task has no reduction, and the process never ends.
% A62
\rcomment{r:e-task-andthen-fail} The continuation $v$ is skipped: a failure
passes over every $\effandThen$ until an $\effonError$ handles it, as rule
\ruleref{r:e-task-onerror-fail} does. Rule \ruleref{r:e-task-onerror-succeed}
is the dual for success.
\rcomment{r:e-task-andthen-inner} Rules \ruleref{r:e-task-andthen-inner} and
\ruleref{r:e-task-onerror-inner} carry both silent and labelled steps out of
the task that is performed first. Every step of a process is a separate
reduction, and another process may take steps between any two of them: a
process may be suspended at any $\effandThen$, and the definition guarantees
no more. In particular, a program cannot rely on a process running without
interruption until it reaches a primitive. % A78
\rcomment{r:e-task-prim} This is the only rule with a label, and it does not
determine its result: the outcome is any member of
$\effOut(\pi,v_1\cdots v_k)$. For $\effspawn$, the process identifier is
unconstrained here; Chapter~\ref{ch:programs} chooses one that names no
earlier process. For the primitives of the platform, the outcome is decided
by the platform, which may also take time to decide it, as for
$\effsleep$. % R9, A83
\end{rulecomments}

The following properties follow from the rules.
\begin{property}
An outcome has no reduction. If $\tk$ is not an outcome, then either every
reduction of $\tk$ is silent, and all have the same right-hand side up to the
choice of a member of a set given by $\APPLY$; or every reduction of $\tk$ is
labelled by an action $\effdone{\pi(v_1\cdots v_k)}{\tk'}$, for one and the
same $\pi$ and $v_1\cdots v_k$; or $\tk$ has no reduction, which happens only
when an application in a premise of rule \ruleref{r:e-task-andthen-succeed}
or \ruleref{r:e-task-onerror-fail} has no result.
\end{property}

\section{Commands and Subscriptions}\label{sec:effects-commands}

A command value, or a subscription value, is a finite tree whose leaves are
requests to managers. A \emph{manager} $\effm$ is one of the following.
\begin{itemize}
\item The \emph{task manager} $\effmtask$. Its leaves $\effleaf(\effmtask,\tk)$
  are commands; it performs $\tk$ in a new process, and when $\tk$ succeeds
  with $v$, the leaf gives a message from $v$.
\item The manager $\effmout\,\effp$ of an outgoing port $\effp$. Its leaves
  $\effleaf(\effmout\,\effp,j)$, with $j\in\BVal$, are commands; each sends
  $j$ to the host, and gives no message.
\item The manager $\effmin(\effp,\tau)$ of an incoming port $\effp$ of payload
  type $\tau$. Its leaves $\effleaf(\effmin(\effp,\tau),v)$ are subscriptions,
  where $v$ is a function value; while the subscription is current, each
  boundary value the host sends on $\effp$ that decodes at $\tau$ gives a message
  through $v$ (Section~\ref{sec:effects-messages}).
\item A manager of the platform, in $\effMgrPlat$. The values its leaves hold,
  whether they are commands or subscriptions, and the messages they give are
  defined by the platform. % A83
\end{itemize}
Since port names are distinct (A84), a port name $\effp$ determines the
payload type $\tau$ of every manager $\effmin(\effp,\tau)$ of a program. A
batch groups commands, and $\effmap(v,\cm)$ applies the function value $v$ to
every message that $\cm$ gives.

A command or subscription value with no $\effbatch$ in it has the form
$\effmap(v_n,\cdots\effmap(v_1,\effleaf(\effm,v))\cdots)$ with $n\geq 0$, and
is called a \emph{leaf command} (or \emph{leaf subscription}); its manager is
$\effm$. The \emph{flattening} $\effFlat(\cm)$ of a command value $\cm$ is the
finite multiset of leaf commands defined by
\[
  \begin{array}{rcl}
  \effFlat(\effbatch(\cm_1\cdots\cm_n)) &=&
    \effFlat(\cm_1)\uplus\cdots\uplus\effFlat(\cm_n),\\
  \effFlat(\effmap(v,\cm)) &=&
    \effmset{\effmap(v,\effc)\mid \effc\in\effFlat(\cm)},\\
  \effFlat(\effleaf(\effm,v)) &=& \effmset{\effleaf(\effm,v)},
  \end{array}
\]
where a leaf command that occurs several times in $\effFlat(\cm)$ gives as
many members of the multiset on the right of the second equation. The
flattening $\effFlat(\sbv)$ of a subscription value is defined by the same
equations. In particular
$\effFlat(\effnone)=\effFlat(\effbatch(\effnone\ \effnone))=\effmset{\,}$.

We have adopted the following convention, the \emph{effect order convention}.
A command determines, through its flattening, which leaf commands it holds
and how many times each, but no order among them. Wherever the leaf
commands of commands are performed (Chapter~\ref{ch:programs}), the leaf
commands of one command for one manager are performed in no specified order,
and the leaf commands for different managers are performed in no specified
order relative to one another. % A66, A67
The same holds for subscriptions. The orders that implementations choose are
recorded in Appendix~\ref{ch:appendix-implementation}.

A program value $\effprog\,r$ is the meaning of a \kw{Program}: $r$ is the
record of the functions that make it up, with at least the fields \kw{init},
\kw{update} and \kw{subscriptions}. The function \kw{Platform.worker} makes a
program value from its record (Figure~\ref{fig:effects-apply}). A platform
may provide further functions whose results are program values, whose
records have further fields, such as \kw{view}; what those fields mean is
defined by the platform. % A82, A83
Chapter~\ref{ch:programs} defines how a program value runs.

\section{The Effect Values of the Basis}\label{sec:effects-apply}

Figure~\ref{fig:effects-apply} gives the arity of each basic value of the
initial basis that makes an effect value, and the result of $\APPLY$ on it
(Section~\ref{sec:dynamic-core-basic}). Each set has exactly one member. One
of these basic values, \kw{Task.command}, is bound to no identifier of the
basis; it is applied only by the texts of \kw{Task.perform} and
\kw{Task.attempt}, as described below
(Appendix~\ref{ch:appendix-dynamic-basis}).

\begin{figure}[htbp]
\centering
\renewcommand{\arraystretch}{1.15}
\small
\begin{tabular}{@{}l@{\quad}c@{\quad}>{$}l<{$}@{}}
\hline
basic value $b$ & $\arity(b)$ & \APPLY(b,v_1\cdots v_n)\\
\hline
\kw{Task.succeed}  & 1 & \APPLY(b,v)=\{\effsucceed\,v\}\\
\kw{Task.fail}     & 1 & \APPLY(b,v)=\{\efffail\,v\}\\
\kw{Task.andThen}  & 2 & \APPLY(b,v,\tk)=\{\effandThen(v,\tk)\}\\
\kw{Task.onError}  & 2 & \APPLY(b,v,\tk)=\{\effonError(v,\tk)\}\\
\kw{Task.command}  & 1 & \APPLY(b,\kw{Perform}\ \tk)=\{\effleaf(\effmtask,\tk)\}\\
\kw{Process.sleep} & 1 & \APPLY(b,v)=\{\effprim(\effsleep,v)\}\\
\kw{Process.spawn} & 1 & \APPLY(b,\tk)=\{\effprim(\effspawn,\tk)\}\\
\kw{Process.kill}  & 1 & \APPLY(b,\pid)=\{\effprim(\effkill,\pid)\}\\
\kw{Cmd.batch}     & 1 & \APPLY(b,[\cm_1,\ldots,\cm_n])=\{\effbatch(\cm_1\cdots\cm_n)\}\\
\kw{Cmd.map}       & 2 & \APPLY(b,v,\cm)=\{\effmap(v,\cm)\}\\
\kw{Sub.batch}     & 1 & \APPLY(b,[\sbv_1,\ldots,\sbv_n])=\{\effbatch(\sbv_1\cdots\sbv_n)\}\\
\kw{Sub.map}       & 2 & \APPLY(b,v,\sbv)=\{\effmap(v,\sbv)\}\\
\kw{Platform.worker} & 1 & \APPLY(b,r)=\{\effprog\,r\}\\
\hline
\end{tabular}
\caption{APPLY on Effect Values}
\label{fig:effects-apply}
\end{figure}

The other functions of these modules that the initial basis provides are
written in Elm in elm/core, and each is defined by its text
(Section~\ref{sec:appendix-dynamic-basis-elm}). They are \kw{Task.map},
\kw{Task.map2} to \kw{Task.map5}, \kw{Task.sequence}, \kw{Task.mapError},
\kw{Task.perform}, \kw{Task.attempt}, \kw{Cmd.none} and \kw{Sub.none}. % A60
Thus \kw{Cmd.none} and \kw{Sub.none} are $\effnone$, the value of
\kw{batch []}. The texts of \kw{Task.perform} and \kw{Task.attempt} apply
\kw{command} to a constructor value $\kw{Perform}\ \tk$, where \kw{Perform}
is the constructor of the type \kw{MyCmd} of the module \kw{Task}, which that
module does not expose. Effect modules are outside Standard Elm, so these two
names are given directly: the texts of \kw{Task} are evaluated in an
environment that binds \kw{Perform} to the constructor \kw{Perform}, of arity
1, and \kw{command} to the basic value \kw{Task.command} of the figure, and
neither is bound to an identifier of the basis, as
Appendix~\ref{ch:appendix-dynamic-basis} states
(Section~\ref{sec:appendix-dynamic-basis-elm}). % A60, A85
Hence \kw{Task.perform $v$ $\tk$} evaluates to
$\effleaf(\effmtask,\tk')$, where $\tk'$ is the value of
\kw{Task.map $v$ $\tk$}, a task that cannot fail and whose success value is
the message. Commands of outgoing ports and subscriptions of incoming ports
are made by the values of the ports (Section~\ref{sec:effects-ports}).

\section{Ports}\label{sec:effects-ports}

A value crosses a port as a boundary value: the port boundary is defined over
the boundary values of BVal, which are the JSON values of JVal together with
those in which byte sequences occur, and host values that are not boundary
values are outside the definition. How a host represents a byte sequence is
outside the definition as well
(Appendix~\ref{ch:appendix-implementation}). % A74, A92, D11
A value $v$ sent on an outgoing port of payload type $\tau$ becomes the
boundary value $\Enc{\tau}(v)$, and a boundary value $j$ arriving at an
incoming port of payload type $\tau$ becomes the value $\Dec{\tau}(j)$, if
that is defined.
The payload types are the types $\tau$ with $\Port(\tau)$
(Section~\ref{sec:static-modules-ports}); aliases have been expanded, and a
payload type has no type variables. Labels are identified with the strings of
their characters, so that a record type's labels are the keys of the objects
that encode its records. The values of the type \kw{Bytes} are the
constructor values $\dbbin\,[o_1,\ldots,o_k]$ of
Appendix~\ref{ch:appendix-dynamic-basis}
(Section~\ref{sec:appendix-dynamic-basis-bytes}), each $o_i$ an Int value
from $0$ to $255$. % A92

\rulegroup{Encoding}{\Enc{\tau}(v)=j}

\defnrule[r:e-enc-int]
  {}
  {\Enc{\tInt}(\sv)=\sv}

\defnrule[r:e-enc-float]
  {\sv\in\effFltFin}
  {\Enc{\tFloat}(\sv)=\sv}

\defnrule[r:e-enc-float-null]
  {\sv\in\effFlt\setminus\effFltFin}
  {\Enc{\tFloat}(\sv)=\effnull}

\defnrule[r:e-enc-true]
  {}
  {\Enc{\tBool}(\kw{True})=\efftrue}

\defnrule[r:e-enc-false]
  {}
  {\Enc{\tBool}(\kw{False})=\efffalse}

\defnrule[r:e-enc-string]
  {}
  {\Enc{\tString}(\sv)=\sv}

\defnrule[r:e-enc-bytes]
  {}
  {\Enc{\tBytes}(\dbbin\,[o_1,\ldots,o_k])=\effbytes(o_1\cdots o_k)}

\defnrule[r:e-enc-unit]
  {}
  {\Enc{\tunit}(())=\effnull}

\defnrule[r:e-enc-nothing]
  {}
  {\Enc{\tau\,\tMaybe}(\kw{Nothing})=\effnull}

\defnrule[r:e-enc-just]
  {\Enc{\tau}(v)=j}
  {\Enc{\tau\,\tMaybe}(\kw{Just}\ v)=j}

\defnrule[r:e-enc-list]
  {\Enc{\tau}(v_1)=j_1\ \ \cdots\ \ \Enc{\tau}(v_n)=j_n}
  {\Enc{\tau\,\tList}([v_1,\ldots,v_n])=\effJarr{j_1,\ldots,j_n}}

\defnrule[r:e-enc-array]
  {\japp{\kw{Array.toList}}{v}{v'}\premsep\Enc{\tau\,\tList}(v')=j}
  {\Enc{\tau\,\tArray}(v)=j}

\defnrule[r:e-enc-tuple]
  {\Enc{\tau_1}(v_1)=j_1\premsep\Enc{\tau_2}(v_2)=j_2\premsep
   \opt{\Enc{\tau_3}(v_3)=j_3}\premsep k=2\opt{{}+1}}
  {\Enc{(\tau_1,\tau_2\opt{,\tau_3})\,\tTuple{k}}((v_1,v_2\opt{,v_3}))
   =\effJarr{j_1,j_2\opt{,j_3}}}

\defnrule[r:e-enc-record]
  {\Dom r=\Dom\varrho\premsep j\in\effStr\finmap\BVal\premsep\Dom j=\Dom\varrho\\
   \Enc{\varrho(\lab)}(r(\lab))=j(\lab)\ \text{for each}\ \lab\in\Dom\varrho}
  {\Enc{\rec{\varrho}}(r)=j}

\defnrule[r:e-enc-value]
  {j\in\JVal}
  {\Enc{\tValue}(j)=j}

\begin{rulecomments}
\rcomment{r:e-enc-int} The integer itself is the JSON number. Whether a host
can hold every integer of $\IntRange$ exactly is not a matter for the
definition (Appendix~\ref{ch:appendix-implementation}). % A73
\rcomment{r:e-enc-float} The sign of a zero is not kept: $-0$ encodes as the
JSON number $0$ (Section~\ref{sec:effects-simple}). % A88
\rcomment{r:e-enc-float-null} JSON has no number for NaN or for an infinity.
Such a float is encoded as $\effnull$, which a \kw{Float} port does not
decode, while a \kw{Maybe Float} port decodes it as \kw{Nothing}. % A74, A88
\rcomment{r:e-enc-bytes} A value of type \kw{Bytes} is sent as its byte
sequence, which is not a JSON value: not a string, and not an array of
numbers. Here each $o_i$ on the left is an Int value and on the right the
same integer. % A92
\rcomment{r:e-enc-just} A \kw{Just} has no encoding of its own. Hence the
encoding of a \kw{Maybe} value cannot tell \kw{Nothing} from a \kw{Just}
whose argument encodes as $\effnull$: in \kw{Maybe (Maybe Int)}, both
\kw{Nothing} and \kw{Just Nothing} encode as $\effnull$, and so do
\kw{Nothing} and \kw{Just ()} in \kw{Maybe ()}.
\rcomment{r:e-enc-array} An array is encoded as the list of its elements.
\kw{Array.toList} is the basic value of
Section~\ref{sec:appendix-dynamic-basis-arrays}, bound to no identifier of the
basis.
\rcomment{r:e-enc-record} The fields become the keys of an object, which has
no order, so the order of fields is no part of what is sent. % A76
\rcomment{r:e-enc-value} A value of type \kw{Json.Encode.Value} is already a
JSON value and is sent as it is. The side condition $j\in\JVal$ holds of every
value of that type, so no byte sequence is ever sent through
\kw{Json.Encode.Value}. % A74, A92
\end{rulecomments}

\rulegroup{Decoding}{\Dec{\tau}(j)=v}

\defnrule[r:e-dec-int]
  {j\in\IntRange}
  {\Dec{\tInt}(j)=j}

\defnrule[r:e-dec-float]
  {j\in\effJNum}
  {\Dec{\tFloat}(j)=\efffloat(j)}

\defnrule[r:e-dec-true]
  {}
  {\Dec{\tBool}(\efftrue)=\kw{True}}

\defnrule[r:e-dec-false]
  {}
  {\Dec{\tBool}(\efffalse)=\kw{False}}

\defnrule[r:e-dec-string]
  {j\in\effStr}
  {\Dec{\tString}(j)=j}

\defnrule[r:e-dec-bytes]
  {}
  {\Dec{\tBytes}(\effbytes(o_1\cdots o_k))=\dbbin\,[o_1,\ldots,o_k]}

\defnrule[r:e-dec-unit]
  {}
  {\Dec{\tunit}(\effnull)=()}

\defnrule[r:e-dec-nothing]
  {}
  {\Dec{\tau\,\tMaybe}(\effnull)=\kw{Nothing}}

\defnrule[r:e-dec-just]
  {j\neq\effnull\premsep\Dec{\tau}(j)=v}
  {\Dec{\tau\,\tMaybe}(j)=\kw{Just}\ v}

\defnrule[r:e-dec-list]
  {\Dec{\tau}(j_1)=v_1\ \ \cdots\ \ \Dec{\tau}(j_n)=v_n}
  {\Dec{\tau\,\tList}(\effJarr{j_1,\ldots,j_n})=[v_1,\ldots,v_n]}

\defnrule[r:e-dec-array]
  {\Dec{\tau\,\tList}(j)=v\premsep\japp{\kw{Array.fromList}}{v}{v'}}
  {\Dec{\tau\,\tArray}(j)=v'}

\defnrule[r:e-dec-tuple]
  {n\geq k\premsep k=2\opt{{}+1}\\
   \Dec{\tau_1}(j_1)=v_1\premsep\Dec{\tau_2}(j_2)=v_2\premsep
   \opt{\Dec{\tau_3}(j_3)=v_3}}
  {\Dec{(\tau_1,\tau_2\opt{,\tau_3})\,\tTuple{k}}(\effJarr{j_1,\ldots,j_n})
   =(v_1,v_2\opt{,v_3})}

\defnrule[r:e-dec-record]
  {j\in\effStr\finmap\BVal\premsep\Dom\varrho\subseteq\Dom j\premsep
   \Dom r=\Dom\varrho\\
   \Dec{\varrho(\lab)}(j(\lab))=r(\lab)\ \text{for each}\ \lab\in\Dom\varrho}
  {\Dec{\rec{\varrho}}(j)=r}

\defnrule[r:e-dec-value]
  {j\in\JVal}
  {\Dec{\tValue}(j)=j}

\begin{rulecomments}
\rcomment{r:e-dec-int} A JSON number is an Int exactly when it is an integer
within $\IntRange$: \kw{1.0} and \kw{-0} decode as \kw{1} and \kw{0}, while
\kw{1.5}, and an integer outside $\IntRange$, do not decode. Since
$\IntRange$ is a parameter of the implementation, so is the set of numbers a
port of type \kw{Int} accepts (Appendix~\ref{ch:appendix-implementation}).
% A73
\rcomment{r:e-dec-float} Every JSON number decodes, integers among them, as
the nearest floating-point number. A number too large for a finite
floating-point number decodes as the infinity of its sign, by the rounding of
floating-point constants. An incoming $-0$ is the number $0$, and decodes as
the float $0$, not $-0$. No boundary value decodes as NaN. % A88
\rcomment{r:e-dec-bytes} A byte sequence, and nothing else, decodes at
\kw{Bytes}: $\Dec{\tBytes}(j)$ is undefined for every JSON value $j$, so a
string or an array of numbers sent to a \kw{Bytes} port is rejected at the
boundary. % A71, A92
\rcomment{r:e-dec-unit} At an incoming port, the unit type accepts
$\effnull$ and nothing else. The flags of a program are decoded otherwise
when their type is the unit type (Chapter~\ref{ch:programs}). % A77
\rcomment{r:e-dec-just} The side condition $j\neq\effnull$ makes
$\effnull$ always decode as \kw{Nothing}, at once, by rule
\ruleref{r:e-dec-nothing}: in \kw{Maybe (Maybe Int)} it decodes as the outer
\kw{Nothing}, and \kw{Just Nothing} is never the result of decoding. A value
other than $\effnull$ that does not decode at $\tau$ does not decode at
\kw{Maybe} $\tau$ either.
\rcomment{r:e-dec-tuple} An array longer than the tuple is accepted, and its
further elements are ignored; a shorter array does not decode.
\rcomment{r:e-dec-record} An object with further keys is accepted, and they
are ignored; an object that lacks the key of a field does not decode. Which
missing field an implementation reports, and how, is outside the
definition. % A76
\rcomment{r:e-dec-value} Every JSON value decodes at \kw{Json.Encode.Value}.
A boundary value in which a byte sequence occurs is not a JSON value, and does
not decode there: it is rejected. % A74, A71, A92
\end{rulecomments}

\begin{figure}[htbp]
\centering
\small
\renewcommand{\arraystretch}{1.15}
\begin{tabular}{@{}>{\raggedright}p{0.24\textwidth}>{\raggedright}p{0.08\textwidth}>{\raggedright}p{0.29\textwidth}>{\raggedright\arraybackslash}p{0.29\textwidth}@{}}
\hline
Elm type & Port? & Sent as (Enc) & Accepted (Dec)\\
\hline
\kw{Int} & yes & the number & a number that is an integer in $\IntRange$\\
\kw{Float} & yes & the number; $\effnull$ for NaN and the infinities & any number, rounded\\
\kw{Bool} & yes & $\efftrue$ or $\efffalse$ & $\efftrue$ or $\efffalse$ only\\
\kw{String} & yes & the string & a string only\\
\kw{Bytes} & yes & a byte sequence & a byte sequence only\\
\kw{()} & yes & $\effnull$ & $\effnull$ only (flags: Chapter~\ref{ch:programs})\\
\kw{Maybe} $\ty$ & if $\ty$ is & $\effnull$ for \kw{Nothing}; $v$ as $\ty$ for \kw{Just} $v$ & $\effnull$ as \kw{Nothing}; anything else as $\ty$, under \kw{Just}\\
\kw{List} $\ty$ & if $\ty$ is & an array & an array whose elements decode as $\ty$\\
\kw{Array} $\ty$ & if $\ty$ is & an array & an array whose elements decode as $\ty$\\
\kw{($\ty_1$, $\ty_2$)}, \kw{($\ty_1$, $\ty_2$, $\ty_3$)} & if each $\ty_i$ is & an array of 2 or 3 elements & an array at least as long; further elements ignored\\
closed record \kw{\kwlb{} $\lab_1$ : $\ty_1$, \ldots\ \kwrb} & if each $\ty_i$ is & an object with exactly the fields as keys & an object with at least those keys; further keys ignored\\
\kw{Json.Encode.}\linebreak\kw{Value} & yes & itself & any JSON value; no byte sequence\\
\hline
\kw{Char} & no & & \\
a function type & no & & \\
a type variable & no & & \\
an extensible record & no & & \\
any other custom type (\kw{Result}, \kw{Order}, \kw{Dict}, \kw{Set}, and those of the program) & no & & \\
any other abstract type (\kw{Cmd}, \kw{Sub}, \kw{Task}, \kw{Process.Id}, \kw{Program}) & no & & \\
\hline
\end{tabular}
\caption{What Can and Cannot Be Sent over a Port}
\label{fig:effects-port-table}
\end{figure}
% A92: Bytes admissible in every payload position, as a byte sequence.

There is no rule for any other pair of a type and a boundary value: a string
does not decode as an Int, nor a number as a Bool, nor an object as a tuple,
nor a byte sequence as a String or a list. The
rules give the following properties.
\begin{property}
For every $\tau$ with $\Port(\tau)$, $\Enc{\tau}$ and $\Dec{\tau}$ are
partial functions: each value has at most one encoding, and each boundary value at
most one decoding, at $\tau$. Every value that elaboration admits as a
payload of type $\tau$ has an encoding at $\tau$. If $\Enc{\tau}(v)=j$, and
$v$ has no component that is NaN, an infinity or a negative zero, and no
component $\kw{Just}\ v'$ whose argument encodes as $\effnull$, then
$\Dec{\tau}(j)=v''$, where $v''$ is $v$ with each array replaced by the value
of \kw{Array.fromList} applied to the list of its elements.
\end{property}

The value of a port is a basic value of arity 1 that the dynamic semantics of
modules binds to the port name (Chapter~\ref{ch:dynamic-modules}). For an
outgoing port $\effp$ of payload type $\tau$, that is the basic value
$\effoutport(\effp,\tau)$, and for an incoming port the basic value
$\effinport(\effp,\tau)$. The class BasVal of
Chapter~\ref{ch:dynamic-core} includes these values, defined here, for every
port name and every payload type. Their meaning is
\[
  \begin{array}{rcl}
  \APPLY(\effoutport(\effp,\tau),v) &=&
    \{\effleaf(\effmout\,\effp,j)\mid\Enc{\tau}(v)=j\},\\
  \APPLY(\effinport(\effp,\tau),v) &=& \{\effleaf(\effmin(\effp,\tau),v)\}.
  \end{array}
\]
By the property above, the first set has exactly one member for every value
an outgoing port is applied to: a value is encoded when the command is made.
% D11
A value sent to an incoming port is decoded when it arrives
(Section~\ref{sec:effects-messages}).

Figure~\ref{fig:effects-port-table} collects what can and cannot be sent over
a port, for each form of type, in agreement with the predicate $\Port$ of
Section~\ref{sec:static-modules-ports}. A type that is not a payload type
makes a port declaration fail to elaborate; no such port exists at run time.
% D11, A92


A type alias stands for the type it expands to. \kw{Bool} and \kw{Maybe} are
custom types, and the only custom types a port admits.

\section{Messages}\label{sec:effects-messages}

A leaf command gives messages from the inputs its manager receives: a task
leaf from the success value of the process that performs its task, an
incoming port leaf from the boundary values that arrive at its port. The rules
below define the message given, through the functions of the $\effmap$s
around the leaf.

\rulegroup{Messages}{\effrecv{\effc}{\effx}{v}}

\defnrule[r:e-msg-task]
  {}
  {\effrecv{\effleaf(\effmtask,\tk)}{v}{v}}

\defnrule[r:e-msg-port]
  {\Dec{\tau}(j)=v'\premsep\japp{v}{v'}{v''}}
  {\effrecv{\effleaf(\effmin(\effp,\tau),v)}{(\effp,j)}{v''}}

\defnrule[r:e-msg-map]
  {\effrecv{\cm/\sbv}{\effx}{v'}\premsep\japp{v}{v'}{v''}}
  {\effrecv{\effmap(v,\cm/\sbv)}{\effx}{v''}}

\begin{rulecomments}
\rcomment{r:e-msg-task} The input $v$ is the value with which the process
performing $\tk$ succeeds. The task of a leaf made by \kw{Task.perform} or
\kw{Task.attempt} cannot fail. For \kw{Task.perform} it is the task given,
mapped, whose error type is \kw{Never}, of which there is no value. For
\kw{Task.attempt} it is the task given with \kw{Task.onError} applied, which
turns a failure into a success whose value is the message made from \kw{Err}.
Either way the task that \kw{Perform} carries has the error type \kw{Never}
(Section~\ref{sec:appendix-static-basis-modenv}). So every such process that
ends gives a message.
\rcomment{r:e-msg-port} The port name of the input must be that of the
manager, so a boundary value arriving at one port gives no message through a
leaf of another. When $\Dec{\tau}(j)$ is undefined, no rule applies, and the
leaf gives no message.
\rcomment{r:e-msg-map} The innermost function is applied first. There is no
rule for a leaf of an outgoing port: such a leaf gives no message, and the
functions of the $\effmap$s around it are never applied. The messages that
the leaves of the managers of the platform give are defined by the platform.
% A83
\end{rulecomments}

A boundary value $j$ that arrives at an incoming port $\effp$ of payload type
$\tau$ is \emph{accepted} if $\Dec{\tau}(j)$ is defined, and otherwise
\emph{rejected}. A rejected value is refused at the boundary: the program
never sees it, and it is unaffected. % A71, D2
When $j$ is accepted and $\sbv$ is the current subscription value of the
program, every member $\effc$ of $\effFlat(\sbv)$ for which
$\effrecv{\effc}{(\effp,j)}{v}$ gives the message $v$, once for each time it
occurs in $\effFlat(\sbv)$. Whether a value is accepted depends on its port
alone, not on the current subscriptions; an accepted value that no current
subscription receives gives no message. When the messages are handled, and
in which order, is defined in Chapter~\ref{ch:programs}.
